\chapter{数与字母：把具体问题变成可以搬运的模型}

\section{一张一百元的购物清单}

开学前一天，妈妈递给你一张一百元的钞票，说：“去买这学期要用的文具，笔记本六元一本，中性笔四元一支，自己搭配，最好刚好把这一百元花完。”

这听起来像一件小事，但站在文具店的货架前，你会发现它其实是个正经的问题：买几本笔记本、几支笔，才能恰好花完一百元？而且妈妈说不定还会追加条件——“笔记本和笔都得买，不能只买一样”。

先凭直觉猜一猜。买十本笔记本是六十元，剩下四十元正好买十支笔——一个答案到手了：$10$ 本、$10$ 支。那还有没有别的买法？$16$ 本笔记本是九十六元，再加一支笔，正好一百。看起来答案不止一个。于是一个更“贪心”的问题冒了出来：

\begin{center}
\textbf{到底有多少种买法？你能保证一种都不漏、一种都不重复吗？}
\end{center}

这个问题会让人卡住。它不像“$10$ 本加 $10$ 支多少钱”那样一步算出结果；它要你找出\emph{所有}答案，还要给出“没有遗漏”的理由。凭感觉报数是不行的——你说有 $8$ 种，同桌说有 $9$ 种，谁对？怎么裁判？

这就是本章的起点。我们会看到：先老老实实用算术去试，试着试着会失控；然后请出本章的主角——\emph{字母}，把“每一道题”变成“一类题”；最后你会得到一个小小的模型，它不只解决这一家文具店的问题，还能原封不动地搬到出租车计费、手机套餐、调配饮料、分压岁钱等一大堆看似无关的场景里。

\section{算术的办法：当“挨个试”开始失控}

最朴素的办法谁都想得出来：挨个试。

笔记本最少 $1$ 本、最多 $16$ 本（$17$ 本就超过一百元了）。对每一种本数，算出剩下的钱，再看看能不能被四整除、商是不是至少为 $1$。于是我们得到一张表：

\begin{center}
\begin{tabular}{ccc}
\toprule
笔记本 $x$（本） & 剩下的钱（元） & 中性笔 $y$（支） \\
\midrule
$2$ & $100-12=88$ & $22$ \\
$4$ & $100-24=76$ & $19$ \\
$6$ & $100-36=64$ & $16$ \\
$8$ & $100-48=52$ & $13$ \\
$10$ & $100-60=40$ & $10$ \\
$12$ & $100-72=28$ & $7$ \\
$14$ & $100-84=16$ & $4$ \\
$16$ & $100-96=4$ & $1$ \\
\bottomrule
\end{tabular}
\end{center}

仔细核对之后，一共 $8$ 种买法。注意一个细节：奇数本笔记本都不可能出现——剩下的钱会是 $94$、$82$ 这样的数，它们除以 $4$ 有余数，买不到整支的笔。我们是靠“逐个数钱”才发现这一点的。

到这里问题似乎解决了。但请看看这个办法的代价：我们做了 $16$ 次减法、$16$ 次除法，还要自己留意“剩下的钱必须刚好分完”。如果妈妈把预算改成一百五十元、两百元，表格就得重做一遍。如果货架上多一种商品——比如两元一张的贴纸，条件变成“三样都要买、刚好花完一百元”——那枚举就不再是一张表，而是几十行甚至上百行：买 $1$ 张贴纸时有哪些买法、买 $2$ 张时又有哪些买法……手工做下去，出错只是时间问题。

算术的失败不在于算错，而在于\emph{每一个新问题都要从头再来}。我们算完一百元，对“一百二十元”毫无积累；数清两种文具，对三种文具毫无办法。我们需要的不是更快的试，而是一种能把“试”本身写下来的办法。

\begin{shiyan}
拿一张纸，自己动手把上面的枚举重做一遍，但把预算改成 $120$ 元（笔记本仍为 $6$ 元、笔仍为 $4$ 元，两样都至少买一件）。数一数你做了多少次除法。然后先别往下看，想一想：这些除法里有没有重复出现的“套路”？如果你必须把同样的套路写给一个不会算术的机器人看，你会怎么写？
\end{shiyan}

\section{先点一次名：我们已经认识的数}

在请字母出场之前，先花一点时间和老朋友们打个招呼——因为从这一章开始，它们会常常出现在同一个式子里。

\emph{自然数} $0,1,2,3,\dots$ 是数数的数。买几本笔记本、分几颗糖，都是自然数的领地。\emph{整数}在自然数之外添上负数 $-1,-2,-3,\dots$，用来表示“反方向”的量：零下三度、欠款五元、电梯地下二层。把整数画在一条直线上，就是我们熟悉的数轴：

\begin{center}
\begin{tikzpicture}[scale=0.85]
\draw[->] (-5.8,0) -- (5.8,0);
\foreach \x in {-5,-4,-3,-2,-1,0,1,2,3,4,5}
  \draw (\x,0.12) -- (\x,-0.12) node[below] {\small $\x$};
\fill[red!60!black] (-3,0) circle (2.2pt) node[above=2pt] {\small 零下 $3$ 度};
\fill[blue!60!black] (2,0) circle (2.2pt) node[above=2pt] {\small 零上 $2$ 度};
\end{tikzpicture}
\end{center}

\emph{分数}（以及小数、百分数这对“近亲”）处理的是“整体的一部分”：半块披萨是 $\frac{1}{2}$，六元的笔记本打七五折就是 $6\times 0.75=4.5$ 元，班级里 $60\%$ 的同学参加了春游。\emph{比例}则研究两个量之间的倍数关系：地图上 $1:10000$ 的比例尺说“图上 $1$ 厘米代表实地 $100$ 米”；把 $12$ 块巧克力按 $2:1$ 分给两人，就是把它切成 $3$ 等份，一人拿两份、一人拿一份。

这些数各有脾气。自然数做除法会“出界”——$7$ 除以 $2$ 在自然数里没答案；整数做除法也一样；分数的世界宽得多，两个分数加减乘除（除数不为零）总还是分数。百分数只是分母固定为 $100$ 的分数，七五折、八折、九五折，说的都是“按原价的百分之多少收钱”。

为什么要在这里点名？因为接下来，当我们写下一个式子时，必须先想清楚：式子里的量允许取哪些数？买文具的数目只能是自然数；温度可以是整数和分数；折扣是百分数。\emph{一个问题的“数的世界”，往往决定了它有没有解、有几个解。}一百元买文具如果允许买“$3.5$ 本笔记本”，买法立刻多得数不清——但文具店不卖半本。这一点我们很快还会用到。

\section{给“不知道”起一个名字：字母与代数式}

回到那张一百元的清单。我们枚举时的每一步其实在做同一件事：假设买了若干本笔记本，算一算，再调整。数学家们很久以前就意识到，与其每次写“假设买了某个数本”，不如干脆给这个“还不知道的数”起一个临时的名字——一个字母。

把笔记本的数量叫作 $x$，中性笔的数量叫作 $y$。花掉的钱就是
\[
6x+4y,
\]
“刚好花完一百元”这个要求就是
\[
6x+4y=100.
\]

请停下来看一眼这行式子，因为它完成了本章第一件大事：\emph{刚才那一整张表格、那十几次除法，被压缩成了一行字}。这行字里没有具体假设买几本，它同时谈论所有可能的买法。预算改成一百二十元？不用重想，只要把式子右边换成 $120$：$6x+4y=120$。文具店换单价？换掉 $6$ 和 $4$ 就行。式子的骨架纹丝不动。

像 $6x+4y$ 这样，由数、字母和运算符号组成的式子，叫作\emph{代数式}。

\begin{shuyu}{代数式}
\textbf{直观解释}：把一道算术题里“会变的那部分”用字母代替后留下的骨架。\\
\textbf{正式定义}：用运算符号（加、减、乘、除、乘方等）把数和表示数的字母连接起来的式子。\\
\textbf{一个例子}：笔记本 $6$ 元一本、笔 $4$ 元一支，买 $x$ 本、$y$ 支的总价是代数式 $6x+4y$；令 $x=10$、$y=10$，它就变回一道普通的算术题 $6\times10+4\times10=100$。
\end{shuyu}

字母在这里扮演的第一个角色是\emph{未知数}：一个具体但暂时不知道的数，我们的工作是通过条件把它“揪出来”。方程 $6x+4y=100$ 里，$x$ 和 $y$ 就是两个未知数，而“$x,y$ 都是正整数”这个来自文具店规矩的限制（回想上一节：数的世界很重要！）将帮我们把答案锁定在有限的几种里。

别小看“起个名字”这一步，数学家们为它摸索了上千年。早期的数学著作里没有符号，一道题要用日常语言写满一整段：“有这样一个量，取它的三倍，再加上二，所得为五十，问这个量是多少。”读的人要先做阅读理解，才能开始算。不同地区的数学家各自发明过各种缩写和记号，彼此看不懂对方的“密码”，争论哪种写法更省事的文章流传了不少。直到字母记号逐渐统一，同一道题才被压缩成 $3x+2=50$ 这七个符号。符号的胜利不在于省墨水，而在于：当条件全部住进式子里，推理就可以对式子直接进行，不再依赖脑内那本账。本章后面所有的变形、求解、推广，都是这场胜利的红利。

用字母写条件还有一个立刻可见的好处：我们能用式子直接\emph{推理}，而不必每次都回到数钱。比如“奇数本笔记本不可能”这件事，用字母说就是：$6x+4y$ 永远是偶数（$6x$ 是偶数，$4y$ 是偶数，偶数加偶数还是偶数），而 $100$ 也是偶数，所以方程两边在“奇偶”上没有矛盾；但 $6x$ 除以 $4$ 的余数随 $x$ 变化——当 $x$ 是奇数时 $6x$ 里有 $6,18,30$ 这样的“$4$ 的倍数加 $2$”，配上 $4y$ 后总和除以 $4$ 余 $2$，而 $100$ 除以 $4$ 余 $0$，矛盾。所以 $x$ 必须是偶数。一句话的推理，代替了表格里八行被划掉的尝试。

\section{天平的守则：等式变形与唯一解}

写下 $6x+4y=100$ 只是开始，我们还要会“解”它。先把术语摆正：

\begin{dingyi}[方程]
含有未知数的等式叫作\emph{方程}。使方程左右两边相等的未知数的值，叫作方程的\emph{解}；求出解的过程叫作\emph{解方程}。
\end{dingyi}

解方程的全部技术，建立在一个古老而朴素的图像上：天平。

一个等式就像一架平衡的天平：左边托盘放着 $6x+4y$，右边托盘放着 $100$，横梁水平。只要不破坏平衡，我们可以对两边做相同的操作——同时添上同样重的砝码，同时拿走，同时翻倍，同时减半。平衡始终保持。

\begin{center}
\begin{tikzpicture}[scale=0.9]
\draw[thick] (-0.7,0) -- (0.7,0) -- (0,1.0) -- cycle;
\draw[thick] (-3.4,1.0) -- (3.4,1.0);
\draw (-3.4,1.0) -- (-3.4,0.35);
\draw (3.4,1.0) -- (3.4,0.35);
\draw (-4.1,0.35) -- (-2.7,0.35);
\draw (2.7,0.35) -- (4.1,0.35);
\node at (-3.4,-0.25) {$3x+2y$};
\node at (3.4,-0.25) {$50$};
\node at (0,1.35) {\small 两边同时减去 $2y$，天平仍然平衡};
\end{tikzpicture}
\end{center}

\begin{dingli}[天平原理]
如果两个式子相等，那么对它们\emph{同时}加上（或减去）同一个数或式子、\emph{同时}乘以同一个数、\emph{同时}除以同一个非零的数，得到的两个式子仍然相等。
\end{dingli}

这条原理听起来平淡，但它是整个方程世界的宪法。比如把 $6x+4y=100$ 两边同时除以 $2$，得到更清爽的 $3x+2y=50$；再把 $2y$ 从两边“拿走”（同时减去 $2y$），得到 $3x=50-2y$；两边除以 $3$：$x=\dfrac{50-2y}{3}$。每一步都只是让天平换一种姿势保持平衡，但走到最后，$x$ 几乎被“单独请了出来”——只要定下一个 $y$，$x$ 立刻被算出来。我们的表格其实就是在 $y=22,19,16,\dots$ 时反复使用这个式子。

天平原理最经典的舞台是一元一次方程——只有一个未知数、未知数只出现一次方的方程。关于它，有一个干净利落的结论：

\begin{dingli}[一元一次方程的解]
设 $a\neq 0$，$b,c$ 是任意已知的数。那么方程 $ax+b=c$ 恰好有一个解
\[
x=\frac{c-b}{a}.
\]
\end{dingli}

\begin{proof}
分两步：先找一个解，再证明没有第二个解。

第一步，用天平原理把 $x$ 解出来。方程两边同时减去 $b$，得 $ax=c-b$；因为 $a\neq0$，两边可以同时除以 $a$，得 $x=\dfrac{c-b}{a}$。把这个值代回原方程检验：$a\cdot\dfrac{c-b}{a}+b=(c-b)+b=c$，确实成立，所以它是解。

第二步，证明唯一。假设 $x_1$ 和 $x_2$ 都是方程的解，即 $ax_1+b=c$ 且 $ax_2+b=c$。两式相减（天平两边减同样的东西），得 $a(x_1-x_2)=0$。由于 $a\neq0$，两个数相乘得零时必有一个是零，所以 $x_1-x_2=0$，即 $x_1=x_2$。所谓“两个解”其实是同一个。证毕。
\end{proof}

请留意证明里的一个小机关：除以 $a$ 之前必须先确认 $a\neq0$。“两边同时除以同一个数”这条守则有个前提——那个数不能是零。除以零没有定义，天平守则在这里有一条明确的边界。这不是吹毛求疵：当 $a=0$ 时方程变成 $0\cdot x+b=c$，若 $b=c$ 则任何 $x$ 都满足，若 $b\neq c$ 则无解——和“恰好一个解”完全不同。\emph{定理的条件和结论一样，都是定理的一部分。}

再看一个数值例子感受一下这条定理的力量。文具店挂出两种优惠：A 方案“满一百减二十五”，B 方案“全场七五折”。买多少钱的东西两种方案一样划算？设原价为 $x$ 元（假设 $x\geq100$，否则 A 不生效），A 方案实付 $x-25$ 元，B 方案实付 $0.75x$ 元。“一样划算”就是 $x-25=0.75x$。两边减 $0.75x$：$0.25x=25$，再除以 $0.25$：$x=100$。原价恰好一百元时两方案打平；超过一百元后 B 方案更便宜（因为 B 的折扣跟着总价涨，A 的减免却固定不变）。百分数、一次方程、天平原理，在一道生活小题里完成了一次合作。

\section{世界不总是刚刚好：不等式}

妈妈的要求如果松一点——“一百元以内随便买，但不能超支”——等式就失效了，因为实付金额 $6x+4y$ 不再等于某个确定的数，而是要满足
\[
6x+4y\leq 100.
\]
这是一个\emph{不等式}。现实中的限制大多长这样：预算不超支、工期不拖延、载重不超标、考试至少及格。等式描述“刚刚好”，不等式描述“够用就行”，后者的地盘其实大得多。

不等式也能像天平一样变形，但守则多了一条警示。同时加上或减去同一个数，方向不变；同时乘以或除以一个\emph{正数}，方向不变；但乘以或除以一个\emph{负数}时，不等号必须掉头。

为什么？看一眼数轴就明白。$2<3$，说的是 $2$ 在 $3$ 的左边。两边同乘 $-1$ 得到 $-2$ 和 $-3$，在数轴上它们关于 $0$ 翻了个面：$-2$ 跑到了 $-3$ 的\emph{右边}，所以 $-2>-3$。乘以负数等于给数轴照镜子，左右当然互换。

\begin{fanli}
警报一：“不等式两边同乘一个数，不等号方向不变。”——错。乘 $-1$ 试试：由 $5<7$ 若得 $-5<-7$ 就闹笑话了，正确的是 $-5>-7$。只有乘\emph{正数}方向才不变，乘负数必须变号，乘零则两边都变成 $0$，不等式直接“死亡”。

警报二：“字母代表的就是那个未知数，解出来就没用了。”——也不对。在 $6x+4y=100$ 里，$x,y$ 可以取 $8$ 组不同的值，字母代表的是“所有候选者”；在下一节你还会看到字母代表“任意一个自然数”。把字母只理解成“藏起来的一个数”，会漏掉它大半的本领。
\end{fanli}

不等式和等式常常联手。回到三种文具的问题：笔记本 $6$ 元、笔 $4$ 元、贴纸 $2$ 元，三样都买、花完一百元。设贴纸买 $z$ 张，那么前两样要满足的正是 $6x+4y=100-2z$，同时 $x\geq1$，$y\geq1$。后两个不等式给 $z$ 划定了活动范围：$100-2z$ 至少要够买一本一笔，即 $100-2z\geq10$，解得 $z\leq45$；而 $z\geq1$。于是 $z$ 只在 $1$ 到 $45$ 之间活动，对每个 $z$ 解一个我们已经会解的二元问题——一个吓人的三元枚举，被不等式修剪成了有边界的、可以分批处理的任务。

\section{字母的另一种工作：装下一整类数}

到目前为止，字母都在替“某个不知道但确定的数”站岗。但字母还有第二份、也许更重要的工作：一次谈论\emph{所有}的数。

看一个古老的规律。把 $1$ 到 $100$ 的自然数全部加起来，和是多少？硬算要算 $99$ 次加法。传说年幼的高斯在课堂上几秒钟就给出了答案 $5050$——这个故事的细节可能经过美化，但其中的想法千真万确，而且用字母一句话就能说清。

\begin{dingli}[连续自然数之和]
对任意正整数 $n$，
\[
1+2+3+\cdots+n=\frac{n(n+1)}{2}.
\]
\end{dingli}

\begin{proof}
把这个和写两遍，一遍正着写、一遍倒着写，上下对齐：
\[
\begin{array}{ccccccc}
1 & + & 2 & + \cdots + & (n-1) & + & n\\
n & + & (n-1) & + \cdots + & 2 & + & 1
\end{array}
\]
每一列上下两个数之和都是 $n+1$：第一列 $1+n$，第二列 $2+(n-1)$，依此类推，最后一列 $n+1$。一共有 $n$ 列，所以两行加起来等于 $n(n+1)$。而这两行是同一个和写了两遍，因此
\[
2\bigl(1+2+\cdots+n\bigr)=n(n+1),
\]
两边同时除以 $2$（天平原理），即得结论。证毕。
\end{proof}

代入 $n=100$：$\dfrac{100\times101}{2}=5050$，与传说一致。但这个定理说的远不止“$1$ 加到 $100$”——它说 $n$ 是 $7$、$1000$、$1000000$ 都对。\emph{一个字母 $n$ 把无穷多道加法题合并成了一道题。}这就是字母的第二份工作：不代表某个未知数，而代表“一整类数里的任意一个”。

这个公式还有一个几何面孔。把 $1+2+\cdots+n$ 想成堆圆点：第一行 $1$ 个、第二行 $2$ 个……第 $n$ 行 $n$ 个，堆成一个三角形。取两个一模一样的“点三角形”，把其中一个倒扣过来，刚好拼成一个 $n$ 行、$n+1$ 列的矩形：

\begin{center}
\begin{tikzpicture}[scale=0.5]
\foreach \j in {1} \fill (\j,-1) circle (3pt);
\foreach \j in {1,2} \fill (\j,-2) circle (3pt);
\foreach \j in {1,2,3} \fill (\j,-3) circle (3pt);
\foreach \j in {1,2,3,4} \fill (\j,-4) circle (3pt);
\foreach \j in {1,2,3,4,5} \fill (\j,-5) circle (3pt);
\foreach \j in {2,3,4,5,6} \draw (\j,-1) circle (3pt);
\foreach \j in {3,4,5,6} \draw (\j,-2) circle (3pt);
\foreach \j in {4,5,6} \draw (\j,-3) circle (3pt);
\foreach \j in {5,6} \draw (\j,-4) circle (3pt);
\foreach \j in {6} \draw (\j,-5) circle (3pt);
\draw[dashed] (0.4,-0.4) rectangle (6.6,-5.6);
\node at (3.5,-6.6) {\small 实心点是 $1+2+3+4+5$，空心点是倒扣的另一份，合起来是 $5\times6$};
\end{tikzpicture}
\end{center}

证明里的“写两遍、上下配对”，在图里就是“拿两份、拼成矩形”。同一个事实，代数说一遍，图形再说一遍——这种双重陈述的习惯，会在本书后面反复出现。

\begin{shuyu}{变量、参数与常量}
\textbf{直观解释}：一场戏里的三种角色。变量是真正在动的演员；参数是这场戏里固定、换一场才换的布景；常量是雷打不动的道具。\\
\textbf{正式定义}：在一个问题中可以取不同值的量叫变量；在同类问题中保持不变、换一类问题才改变的量叫参数；永远不变的数叫常量。\\
\textbf{一个例子}：在 $6x+4y=100$ 中，$x,y$ 是变量，$6,4,100$ 是常量。把问题推广成 $px+qy=m$（单价 $p,q$、预算 $m$），$x,y$ 仍是变量，而 $p,q,m$ 成了参数——换一家文具店、换一个预算，只是换一组参数，模型原封不动。
\end{shuyu}

分清这三种角色，是从“做一道题”跨到“造一个模型”的关键一步。写下 $6x+4y=100$，你只解决了妈妈的这一张清单；写下 $px+qy=m$，你同时解决了\emph{所有}单价、\emph{所有}预算下的同类问题。参数越多，模型搬得越远。

\section{换一个视角：模型是可以搬运的}

现在回头看，本章其实只干了一件事：把具体问题翻译成由字母、等式和不等式组成的模型。值得强调的是，同一个模型会在完全不同的场景里反复现身——这正是代数“一本万利”的地方。

出租车计费：起步价 $10$ 元，之后每公里 $2$ 元。乘 $x$ 公里（$x\geq3$，假设 $3$ 公里内算起步价）的车费是 $2(x-3)+10$，即 $2x+4$ 元。手里只有 $30$ 元，最远能坐多远？解 $2x+4\leq30$，得 $x\leq13$。手机套餐：月租 $19$ 元含 $10$ GB 流量，超出部分每 GB 收 $3$ 元，月消费上限 $50$ 元，最多能用多少流量？设超出 $x$ GB，则 $19+3x\leq50$，得 $x\leq\frac{31}{3}$，即最多超出 $10$ GB。调饮料：浓缩汁与水按 $1:5$ 调配，要配 $1.2$ 升，浓缩汁要多少？设浓缩汁 $x$ 升，$x+5x=1.2$，$x=0.2$ 升。

看出共同点了吗？车费、流量、饮料，剥掉外壳之后都是“形如 $ax+b$ 的量，配上一个等式或不等式”。你在文具店里学会的那套天平守则，在出租车里、在套餐里、在厨房里照样管用——一个字都不用改。

这正是本章最想让你带走的一句话：\emph{代数不是把同一道题算得更快，而是一次解决许多相似的问题}。枚举的努力随问题的规模水涨船高——预算翻倍，表格就翻倍；而模型的努力只花一次，往后的每一个新问题都变成“换一组参数、解一遍方程”的例行公事。算术回答“这一题的答案是多少”，代数回答“这一类题的答案长什么样”。从“一个”到“一类”，就是本章标题里“可以搬运”四个字的全部含义。

\begin{lianjie}
这种“翻译”能力是现代世界的底层技术之一。导航软件把道路网翻译成变量之间的限制，替你解一个庞大的“怎样最快到达”的问题；电商后台把库存、优惠券、满减规则翻译成成百上千个像 $6x+4y\leq100$ 这样的式子，自动判断每个订单能不能成交、该收多少钱；你写程序时声明的每一个“变量”，正是本章字母的电子化身。反过来，本书后面每一章都在给这个翻译器增加新语种：第 2 章加上“变化”，第 3、4 章加上“形状”，再往后还有“随机”和“无限”。
\end{lianjie}

也要诚实地说出模型的边界。把行程压成 $2x+4$ 时，我们忽略了堵车和绕路；把优惠比成 $x-25$ 与 $0.75x$ 时，我们假设你真会把钱花到一百元以上。\emph{模型是对现实的翻译，而翻译永远有取舍。}好的模型抓住问题的主要结构、舍掉无关的细节；判断取舍是否得当，是建模者和数学家共同的功课，本书第 2 章和第 18 章还会回到这个话题。

\begin{pandeng}
继续向前的话，本章的每条线都通向一个更大的世界。要求未知数取整数的方程（如 $6x+4y=100$ 的“文具店版本”）叫\emph{丢番图方程}，是数论的经典猎场，本书第 13 章会专门拜访；一堆不等式围出可行区域、再在其中找最优方案，是\emph{线性规划}的用武之地，运筹学和经济学每天靠它调度真实的货物与资金；“两边做同样操作保持平衡”这条守则本身，在大学的\emph{抽象代数}里被提炼成群与域的公理——数学家们问：还有哪些世界里也存在着自己的“天平”？关键词：丢番图方程、线性规划、公理化。
\end{pandeng}

\begin{toolbox}
本章新增的工具，请收好：
\begin{itemize}
\item \textbf{设字母}：给不知道的数（未知数）或一整类数（任意数）起名字，是建模的第一步。
\item \textbf{代数式}：把“会变的部分”抽出来，留下可以反复使用的骨架，如 $6x+4y$。
\item \textbf{天平原理}：等式两边同加、同减、同乘、同除（除数非零）保持相等——解方程的全部依据。
\item \textbf{不等式变形}：乘除负数要变号；等式管“刚好”，不等式管“范围”，两者常联手。
\item \textbf{区分角色}：变量在动、参数随问题而换、常量不动；参数化是让模型可搬运的关键。
\item \textbf{先问数的世界}：未知数允许取自然数、整数还是任意数？这决定解有几个、甚至有没有。
\end{itemize}
\end{toolbox}

\section*{思考题}

\textbf{观察题}

\begin{sikao}
\item 再看一遍一百元清单的解：$(x,y)=(2,22),(4,19),(6,16),\dots$。每换一种买法，$x$ 增加 $2$，$y$ 减少 $3$。为什么恰好是“加 $2$ 减 $3$”，而不是别的搭配？（提示：假设从一种买法改成另一种，$6x+4y$ 的总额必须不变；让 $x$ 增加某个数，算算 $y$ 该怎么配合。）
\item 在数轴上标出满足 $6x+4y\leq100$（且 $x\geq1$，$y\geq1$）的几组 $(x,y)$——把 $x$ 画在横的方向、$y$ 画在竖的方向，一个买法对应一个点。观察这些点的排列，你有什么发现？（提示：先描 $6x+4y=100$ 对应的点，再看看其他点落在它的哪一侧。）
\end{sikao}

\textbf{证明题}

\begin{sikao}
\item 证明：任意两个相邻整数之和必为奇数。（提示：把较小的那个整数叫作 $n$，把另一个用 $n$ 表示出来再相加。）
\item 证明：$1+3+5+\cdots+(2n-1)=n^2$，即从 $1$ 开始的前 $n$ 个奇数之和是一个完全平方数。（提示：模仿本章“写两遍配对”的办法；或者画点阵——每次在角上添一个“拐弯”形的点层。）
\end{sikao}

\textbf{挑战题}

\begin{sikao}
\item 三种文具问题：$6x+4y+2z=100$，$x,y,z$ 都是正整数。设计一套不遗漏、不重复的计数流程，把买法的总数数清楚。（提示：先固定 $z$，把问题化回本章已解决的二元情形；注意 $6x+4y=100-2z$ 要求右边满足什么条件才有解。）
\item 用 $40$ 米长的围栏靠墙围一块矩形菜地（墙算作一条边，围栏只围三条边），长和宽都取整数米。怎样围面积最大？（提示：设与墙平行的边为 $x$ 米，另两条边各为 $y$ 米，先用围栏长度写出 $x,y$ 的关系，再把面积单独用 $x$ 表示出来，列表试几个值；猜猜不靠墙的版本答案会怎么变。）
\end{sikao}

\section*{通往下一章的门}

挑战题的最后一问里，藏着一个本章工具回答不彻底的问题。围栏总长固定时，$x$ 和 $y$ 被绑在一起：$x$ 每变一点，$y$ 被迫跟着变，面积 $S$ 也跟着变。我们用“列表试几个值”找到了最大的面积，但心里不踏实：整数之外呢？$x$ 取 $13.5$ 米会不会更好？我们试了几个点，却声称看见了全部。

本章的办法是盯着一个一个的数；而这里的真正主角是\emph{变化本身}——$S$ 怎样随着 $x$ 连续地起伏。数学家们为这种问题准备了一件新武器：把关系画到坐标纸上。横的方向表示 $x$，竖的方向表示 $S$，每一个可能的围法变成纸上的一个点，所有点连起来，就是这条关系的“画像”。画像一出，最高点在哪里、从哪边爬上去、往哪边走下坡路，一目了然——不再需要猜测，也不再怕漏掉 $13.5$ 这样的中间值。

这条“画像”就是\emph{函数}的图像。把同一个关系从式子翻译成图形，我们会得到一整套新问题：图像为什么长成抛物线的样子？斜的陡缓说明什么？两条图像相交意味着什么？第 2 章，我们从解一个方程，走向研究所有答案排成的形状。
